% FSE 2027 anonymous research-track submission.
% Revised against the supplied current-evidence report and author corrections.
% This manuscript revision adds no proof targets or new experiment results.
\documentclass[acmsmall,screen,review,anonymous]{acmart}
\setcopyright{none}
\acmDOI{}
\acmISBN{}
\acmYear{2027}
\copyrightyear{2027}
\acmConference[FSE 2027]{ACM International Conference on the Foundations of Software Engineering}{July 12--16, 2027}{Shenzhen, China}

\usepackage{amsmath}
\usepackage{array}
\usepackage{booktabs}
\usepackage{tabularx}
\usepackage{xspace}
\usepackage{graphicx}
\usepackage{url}
\usepackage{tikz}
\usetikzlibrary{arrows.meta,positioning}

\setlength{\emergencystretch}{1.5em}

\newcommand{\Unknown}{\textsc{Unknown}\xspace}
\newcommand{\Commit}{\textsc{Commit}\xspace}
\newcommand{\Fail}{\textsc{Fail}\xspace}
\newcommand{\Retry}{\textsc{Retry}\xspace}
\newcommand{\ReadOnly}{\textsc{ReadOnly}\xspace}
\newcommand{\Idempotent}{\textsc{Idempotent}\xspace}
\newcommand{\Deduplicated}{\textsc{Deduplicated}\xspace}
\newcommand{\Uncontrolled}{\textsc{Uncontrolled}\xspace}
\newcolumntype{Y}{>{\raggedright\arraybackslash}X}
% acmart supplies theorem, proposition, corollary, and proof environments.
\providecommand{\Description}[1]{}

\title[Verified Request-Level Outcomes]{Verified Request-Level Outcomes for Recovering Agent Tool Calls}

% Double-anonymous submission: author and affiliation fields are intentionally
% omitted.

\begin{document}

\begin{abstract}
Agent runtimes use tool outcomes to decide what to do next. After a crash,
a logical request may span several physical attempts, each leaving different
evidence about its effect. We present a verified recovery protocol that
derives request-level \Commit, \Fail, and \Unknown outcomes from this
retained evidence. Its central invariant connects every physical invocation
to durable authorization and attempt ancestry; reporting an idempotent
request as \Fail requires matching failure evidence for every invocation.
Broker--Journal--typed-WAL refinement preserves this invariant through crash
and replay. Three service-adapter laws connect the evidence to external
effects, with operational instances discharging the laws under explicit
interference assumptions. The machine-checked per-execution composition
supports authorization-to-effect consistency, a paper-level specification
of sound outcomes across executions sharing a pre-terminal observation.
A verified typed-record kernel gates selected Rust paths, and all 21
injected-crash cases satisfy their conformance assertions. Paired Redis
broker crashes produce identical pre-terminal WAL evidence and different
service states. The classifier reports \Unknown in both; a targeted
failure-coverage ablation produces an unsound \Fail in the execution with an effect.
Together, the proofs and experiments establish how a serialized recovering
broker can turn retained attempt evidence into justified request-level outcomes.
\end{abstract}

\ccsdesc[500]{Software and its engineering~Formal software verification}
\ccsdesc[300]{Computer systems organization~Reliability}
\ccsdesc[300]{Security and privacy~Authorization}
\keywords{AI agents, tool calls, crash recovery, authorization, effects,
formal verification, refinement}

\maketitle
\section{Introduction}
\label{sec:introduction}

AI agents invoke external tools to deploy software, publish artifacts, modify
infrastructure, and communicate with users. A runtime authorizes a request
and returns an outcome that guides the
planner's next step. We study how a recovering runtime can give that
outcome a precise request-level meaning: \Commit
for an established effect, \Fail for an established absence of effect, and
\Unknown when the retained evidence leaves the effect unresolved.

A crash can separate the service effect from the evidence available to
the runtime.  Consider an agent authorized to execute
\texttt{set\_deployment(v2)}.  The runtime records the request and invokes the
deployment service.  The service changes the active version to \texttt{v2}, but
the runtime crashes before recording the successful reply.  After recovery,
it retries the same request.  By then, its service credential has expired, so
the second attempt returns \texttt{401 Unauthorized} without changing the
deployment.  Figure~\ref{fig:motivating-example} summarizes the execution
and the two external histories compatible with the retained evidence.

\begin{figure}[t]
  \centering
  \includegraphics[width=\linewidth]{figures/motivating-example.pdf}
  \caption{A lost reply leaves zero-effect and one-effect executions
  compatible with the same request-scoped durable observation.}
  \label{fig:motivating-example}
  \Description{A lost first reply and a failed retry leave identical durable evidence compatible with both zero-effect and one-effect executions.}
\end{figure}

A last-response-wins runtime would report \Fail, although the first attempt
may already have changed the deployment. The retained observation also
admits an execution in which the first attempt never reached the service.
These two branches, shown in Figure~\ref{fig:motivating-example}, require
\Unknown: the latest response describes one attempt, while the caller needs
an outcome for the entire logical request. Idempotency makes repeated
updates compatible with the specified effect; resolving the request still
requires evidence about the earlier attempt.

The broker records an attempt-specific Start before invocation. Every
invocation therefore has durable ancestry, while a crash may leave a Start
whose call never occurred. Our central invariant connects retained failure
records to all physical invocations before an idempotent request can be
reported as \Fail. Section~\ref{sec:architecture} develops this argument.

We specify sound outcomes using \emph{authorization-to-effect consistency}
(A2E): a terminal effect claim must hold in every modeled execution sharing
the pre-terminal observation. The proof combines a recovery invariant with
service laws that interpret the evidence. This outcome interface matters
for agents because a reported failure may prompt the planner to issue a
replacement request after the original change has already occurred.

Durable workflows and retained service decisions provide established recovery
mechanisms~\cite{temporal2026activity,lee2015rifl}. Our proof connects a
broker's retained evidence to the physical invocations and service laws
that justify its request-level outcomes.

This paper makes three contributions:
\begin{itemize}
  \item \textbf{Verified request-level outcome reporting.}
  We give \Commit, \Fail, and \Unknown explicit effect meanings and
  establish the evidence invariant that supports them. Every invocation
  retains durable authorization ancestry, and an idempotent \Fail covers
  every physical invocation with matching failure evidence. Journal and
  typed-WAL simulations preserve this invariant through crashes and replay.
  \item \textbf{A compositional proof interface for service adapters.}
  We separate the broker's causal evidence from its service-specific
  effect interpretation through three adapter laws. EnsureMember and
  keyed-decision instances discharge distinct semantic obligations and
  reuse the same WAL terminal-refinement theorem. Covered service
  couplings transport the result to modeled protected state.
  \item \textbf{An executable connection and evaluated artifact.}
  A verified typed-record kernel gates selected Rust append and recovery
  paths. A 21-case crash campaign checks conformance, while paired Redis
  process crashes and a targeted ablation test the failure-coverage rule.
  Proof dependencies characterize structural reuse, and microbenchmarks
  measure local mediation cost.
\end{itemize}

\section{System Model}
\label{sec:system-model}
The broker authorizes tool requests, records attempt evidence, invokes
external services, and exports request-level decisions to the planner.
Recovery uses retained evidence and the applicable service contract.

\subsection{Requests and Attempts}

A planner submits a \emph{logical request} $r$, identifying a canonical
tool operation and its parameters. The broker executes the request
under an associated authorization. A configuration $C$ fixes the
request identity, capability constraints, initial budget, attempt bound,
and recovery profile, together with a stable key when required by the
service.

Authorization is checked at admission and retained for later attempts of
that exact request. Revocation prevents new authorizations; admitted
requests remain subject to their profile, attempt limit, and recovery state.
A remote credential rejection is a service response. The authorization
claim establishes the provenance of admitted invocations; credential
validity and issuance-policy correctness are separate premises.

A \emph{physical attempt} $(r,a)$ is one invocation of $r$. Multiple
attempts preserve the canonical operation and parameters and remain
subject to authorization and recovery constraints; changed parameters
constitute a new request. Attempt outcomes support one terminal decision
for $r$. The recovery profile identifies the service guarantees defined
in Section~\ref{sec:service-laws}; each instance must discharge its adapter obligations.

\subsection{Execution and Durable Evidence}

The external service owns the affected state and shares no transaction
with the broker. Invocation, service effect, response delivery, and
outcome persistence are distinct events.

For an execution $\tau$, let $J(\tau)$ be the sequence of records whose
persistence has been acknowledged to the broker. We call this sequence the
\emph{durable history}. It contains the retained authorization and attempt
records, validated outcomes, and terminal decisions. A response enters
$J(\tau)$ after validation and acknowledged persistence.

For request $r$, let
\[
  J_r(\tau)=\mathsf{Obs}_r(J(\tau))
\]
be its request-scoped \emph{durable observation}, retaining authorization,
attempt, outcome, and terminal records. \emph{Durable evidence} consists
of facts derived from this observation. \emph{Request-wide evidence}
connects these records to the relevant physical attempts: provenance
identifies the authorized request and attempt, and complete failure
coverage accounts for every invocation when an idempotent request selects \Fail.

Let $E_r(\tau)$ denote the external semantic outcome of request $r$
under its service contract. Durable history need not determine this
outcome. In particular, the model admits executions such that
\[
  J_r(\tau_1)=J_r(\tau_2)
  \qquad\text{but}\qquad
  E_r(\tau_1)\ne E_r(\tau_2).
\]
The lost-reply example exhibits this relation.
For the formal recovery model, crashes are fail-stop: they discard volatile
broker state but do not undo service effects. The same durable observation gap
can also arise from lost, delayed, or duplicated responses. Responses may be
delivered after recovery, and recovery may itself crash.
We consider serialized broker execution. The abstract persistence
interface preserves acknowledged records; Section~\ref{sec:architecture} establishes the
corresponding preservation result through the modeled storage layers.

\subsection{Service Contracts and Recovery Decisions}

An \emph{adapter} comprises an API wrapper and a semantic contract
relating accepted responses and repeated attempts to external state
under explicit environmental assumptions.

We write the recovery classifier's interface as
\[
  \mathsf{ED}_C(J,r)
  \in
  \{\Commit,\Fail,\Unknown,\Retry\}.
\]
Here, $C$ supplies the request configuration and recovery profile;
$J$ supplies retained evidence, including the number of started attempts.
The attempt limit bounds invocations; the capability budget bounds
authorization consumption. The service contract supplies the effect
interpretation of the evidence.

\section{Authorization-to-Effect Consistency}
\label{sec:contract}

A2E specifies the effect safety of request-level terminal decisions.

\subsection{Execution Semantics and Terminal Specification}

Fix the request configuration $C$ and adapter $A$. Let
$\mathsf{Exec}_{C,A}$ denote finite executions of the modeled request
interface paired with external runs satisfying $\mathsf{Rely}_A$.
Membership is determined by the operational rules, record-to-event
correspondence, and service rely condition, independently of the terminal
effect conclusion. The rely condition specifies admitted service behavior
and environmental interference.

Use $J_r(\tau)$ and $E_r(\tau)$ from Section~\ref{sec:system-model}. For terminal
soundness, let $\mathsf{DecExec}_{C,A}(r)$ contain the executions in
$\mathsf{Exec}_{C,A}$ cut immediately after acknowledgement of the
first terminal record for $r$, allowing any terminal kind. Let
$J^{-}_r(\tau)$ be the full journal prefix strictly preceding that
record, and define the evidence projection
\[
  H_r(\tau)=\mathsf{Obs}_r(J^{-}_r(\tau)).
\]
Unlike $J_r$, $H_r$ omits the terminal being justified. Write $H$ for
a value of this request-scoped evidence projection. Define
\[
\begin{aligned}
\mathsf{Worlds}_{C,A}(H,r)
  &=\{\tau\in\mathsf{DecExec}_{C,A}(r)\mid H_r(\tau)=H\},\\
\mathsf{Compat}_{C,A}(H,r)
  &=\{E_r(\tau)\mid\tau\in\mathsf{Worlds}_{C,A}(H,r)\}.
\end{aligned}
\]
The class contains all admitted executions sharing $H$, independently of
the candidate decision $d$. Effects are evaluated at the first terminal
acknowledgement; subsequent state changes require separate reasoning.

For a terminal decision $d$, let $\llbracket d\rrbracket_{r,A}$ denote
the external outcomes permitted by its effect interpretation. For \Commit,
the returned value is a component of $d$, suppressed below, and
$\mathsf{resultSpec}_r$ is evaluated on that same value:
\[
\begin{aligned}
\llbracket\Commit\rrbracket_{r,A}
  &= \{E \mid E \models \mathsf{resultSpec}_r
      \land \mathsf{One}_r\}\quad\text{(mutating)},\\
\llbracket\Commit\rrbracket_{r,A}
  &= \{E \mid E \models \mathsf{resultSpec}_r
      \land \mathsf{Zero}_r\}\quad\text{(read-only)},\\
\llbracket\Fail\rrbracket_{r,A}
  &= \{E \mid E \models \mathsf{Zero}_r\},\\
  \llbracket\Unknown\rrbracket_{r,A}
   &= \mathcal{E}_{r,A}.
\end{aligned}
\]
Here $\mathcal{E}_{r,A}$ is the adapter-defined space of external outcomes
for request $r$: \Unknown leaves the effect unresolved. The separate Bound
law constrains admitted runs to $\mathsf{Zero}_r\lor\mathsf{One}_r$.
We write
\[
  H,r,C,A\ \vdash_{\mathrm{A2E}}\ d
\qquad\text{iff}\qquad
  \mathsf{Compat}_{C,A}(H,r)\subseteq\llbracket d\rrbracket_{r,A}.
\]
A2E specifies terminal effect safety; \Retry is a continuation action.
Section~\ref{sec:architecture} derives this inclusion for realized observations, whose execution classes are nonempty.

\subsection{The Decision Contract}

The full decision contract combines effect safety with protocol legality
and request-wide evidence. \Commit requires matching successful-result
evidence; \Fail requires failure evidence resolving the whole request,
including earlier attempts. This distinguishes a failed request from a
successful read-only operation, both of which satisfy $\mathsf{Zero}_r$.
Already-satisfied updates may also satisfy both $\mathsf{Zero}_r$ and
$\mathsf{One}_r$. The following proposition identifies unavoidable ambiguity.

\begin{proposition}[Indistinguishable executions]
\label{prop:indistinguishable}
Consider a mutating request and an observation $H$. Suppose
$\tau_0,\tau_1\in\mathsf{Worlds}_{C,A}(H,r)$ satisfy
\[
\begin{aligned}
E_r(\tau_0)&\models\mathsf{Zero}_r\land\neg\mathsf{One}_r,\\
E_r(\tau_1)&\models\mathsf{One}_r\land\neg\mathsf{Zero}_r.
\end{aligned}
\]
A classifier whose inputs are only $H$, $r$, $C$, and $A$ cannot
soundly return \Commit or \Fail on this observation. If it exports a
terminal from the present vocabulary, that terminal must be \Unknown.
\end{proposition}
\begin{proof}
Identical inputs select the same decision, but \Commit is false in
$\tau_0$ and \Fail is false in $\tau_1$. Thus only \Unknown is a
sound terminal; \Retry may gather evidence but is not a terminal claim.
\end{proof}

The explicit negations distinguish the two outcomes even when the effect
relations overlap for already-satisfied updates. This paper-level observation
argument applies to the lost-reply example when the requested change is
initially absent.

These meanings require evidence with both valid provenance and
persistence. The recovery protocol has three additional historical
obligations:
\begin{enumerate}
  \item \textbf{Authorization ancestry.}
  Every physical invocation is traceable to the exact logical
  request and acknowledged authorization permitting that invocation.
  Capability, budget, and attempt constraints are respected.
  A terminal decision relies only on evidence for that same request.

  \item \textbf{Evidence preservation.}
  Recovery preserves the contents and provenance of acknowledged
  evidence. Accepted outcomes identify their request and attempt;
  an absent outcome is not replaced by an assumed success or failure.

  \item \textbf{Terminal integrity.}
  Each logical request has at most one durable terminal record.
  A \Commit cites the exact successful outcome supporting its
  returned value, and recovery preserves previously established
  terminal decisions.
\end{enumerate}

Sections~\ref{sec:service-laws} and~\ref{sec:architecture} connect these historical obligations to service laws
and establish their preservation through execution and storage refinement.

\subsection{Evidence-Based Recovery Decisions}

For a request in the replay-derived Armed phase without a terminal
decision, recovery examines the retained attempt history.
Its next decision belongs to
\[
\begin{aligned}
  \mathcal{T} &= \{\Commit,\Fail,\Unknown\},\\
  \mathcal{D} &= \mathcal{T}\cup\{\Retry\}.
\end{aligned}
\]
Elements of $\mathcal{T}$ terminate the request. \Retry starts the
first attempt or continues with another permitted attempt. Outside
the Armed phase the mechanized classifier returns Stable; that
protocol-control result is omitted from the decision interface here.

The classifier selects the first matching guard below. Result values,
supporting attempt identifiers, and reason codes are suppressed.

\[
\mathsf{ED}_C(J,r)=
\begin{cases}
  \Commit
    & \text{valid durable success},\\
  \Fail
    & \text{conclusive failure},\\
  \Unknown
    & \text{profile-specific non-conclusive failure},\\
  \Retry
    & \text{eligible attempt below the configured limit},\\
  \Unknown
    & \text{otherwise}.
\end{cases}
\]
\paragraph{Success evidence.}
A valid durable success is an acknowledged successful outcome
whose request identity, attempt provenance, and returned value
satisfy the configured checks. This evidence takes precedence
over failure evidence. The service contract connects the accepted
success to the request's specified effect semantics.

\paragraph{Failure evidence.}
Conclusive failure evidence resolves the whole request under its service
contract, even if a local attempt outcome is missing. For an idempotent
request, a Failure is non-conclusive whenever an earlier started attempt
lacks failure evidence. This branch selects \Unknown even with attempts
remaining. Ambiguous and InvalidResult may instead permit \Retry under the
profile and attempt-limit guards. A deduplicated failure uses an
authoritative keyed decision; transport and credential rejections qualify
only when the service contract gives them request-wide failure meaning.

If neither success nor failure evidence selects a terminal, let $n$ be the
number of started attempts and $m$ the configured maximum. The retry guard is
$n<m\land(n=0\lor\mathsf{class}(r)\ne\Uncontrolled)$.
The $n=0$ case permits the first invocation; later invocations require
a repetition-permitting profile. These attempt bounds are distinct
from the authorization budget.

When no guard matches, the classifier returns \Unknown for residual uncertainty,
including exhausted attempt allowance and missing outcomes. In the
deployment example, the unresolved first attempt initially permits
\Retry. Once the retry records a failure, the non-conclusive-failure
branch selects \Unknown.

\section{Service Laws and Instantiations}
\label{sec:service-laws}

The protocol connects retained records to physical responses; adapter
laws connect those responses to external effects.

\subsection{Semantic Obligations}

An adapter supplies relations over protected external runs for request
$r$: $\mathsf{Zero}_r$ represents absence of a protected effect attributable
to $r$, and $\mathsf{One}_r$ represents the requested abstract effect.
These state relations may overlap for already-satisfied updates, and repeated
insertions can satisfy $\mathsf{One}_r$. Bound requires the adapter to
represent partial mutations and secondary effects, such as billing or
notifications, or exclude them through its rely condition. This obligation
also applies to a single invocation.

The adapter also supplies a result predicate $\mathsf{resultSpec}$
and a rely condition $\mathsf{Rely}$. The latter specifies admitted
service behavior and environmental interference. All effect
guarantees below are relative to this condition.

\paragraph{Records and physical histories.}
Let $\eta_r(\tau)$ be the request-local projection of physical invocations
and deliveries, including deliveries whose results were lost before
persistence. The proof relates this projection to the durable observation
of the same execution.

Two mechanized predicates express request-wide evidence. With configuration
arguments suppressed, $\mathsf{OutcomeEvidence}(J,r,\eta,d)$ establishes
provenance for the selected terminal: an outcome reference backed by a
matching delivery, or an \Unknown reason with a durable anchor.
$\mathsf{BrokerOutcomeCompatible}(J,r,\eta,d)$ supplies the profile-specific
physical-history premises required by the service laws.

For \Commit, compatibility requires the selected successful delivery
for the exact request, attempt, and returned value. For \Fail in the
idempotent profile, it requires
\[
  \forall a\in\mathsf{Invoked}_r(\eta).\;
  \mathsf{Delivery}(\eta,r,a)=\mathsf{Failure}.
\]
The broker derives this coverage from durable failure records and
invocation-to-record correspondence. Deduplicated compatibility additionally
uses retained-decision consistency from the rely condition; uncontrolled
compatibility uses the broker invocation bound. Section~\ref{sec:architecture} establishes these
predicates for the selected terminal.

\paragraph{Primitive laws.}
The first two laws interpret compatible observations. The third
bounds the effect of every execution admitted by the rely condition,
including executions with unresolved outcomes. With common
configuration, request, physical-history, and external-run arguments
omitted, and with $\mathsf{CommitCompatible}$ and
$\mathsf{FailCompatible}$ denoting the corresponding branches of
$\mathsf{BrokerOutcomeCompatible}$, they are:

\begin{align}
  &\mathsf{Rely}\land\mathsf{CommitCompatible}
    \Rightarrow
    \mathsf{resultSpec}\land
    \begin{cases}
      \mathsf{Zero}_r, & \text{read-only},\\
      \mathsf{One}_r, & \text{mutating},
    \end{cases}
    && \tag{Success}\\[0.5ex]
  &\mathsf{Rely}\land\mathsf{FailCompatible}
    \Rightarrow\mathsf{Zero}_r,
    && \tag{Failure}\\[0.5ex]
  &\mathsf{Rely}
    \Rightarrow\mathsf{Zero}_r\lor\mathsf{One}_r.
    && \tag{Bound}
\end{align}

The operational instances below discharge these laws by deriving their rely
conditions from enabled transitions and inductive invariants.

\begin{table}[t]
  \caption{Proof responsibilities and their semantic boundaries.}
  \label{tab:proof-responsibilities}
  \centering
  \small
  \begin{tabularx}{\columnwidth}{@{}lYY@{}}
    \toprule
    \textbf{Obligation} & \textbf{Established by} & \textbf{Conclusion} \\
    \midrule
    History preservation & Broker invariants and Journal/WAL simulations &
      Retained identity, ancestry, and acknowledgement prefix \\
    Causal evidence & Record-to-invocation and delivery proofs (T6-E0) &
      A selected record has matching physical provenance \\
    Profile compatibility & Branch guards and relevant rely conditions
      (T6-C0/S0) & Physical history satisfies the selected law's premises \\
    Effect interpretation & Adapter-specific Success, Failure, and Bound
      proofs (T6-D0/A0 and instances) & Result and protected-state relation
      under the declared service contract \\
    Service transport & Covered operational/service coupling and closed
      interface & The relation holds for the modeled protected service \\
    \bottomrule
  \end{tabularx}
\end{table}

\subsection{Idempotent Operations}

State idempotence, $f_r(f_r(x))=f_r(x)$, bounds repeated state updates.
The adapter establishes Bound over the complete protected effect surface
and Success for accepted results. An exported \Fail decision requires
the invocation-wide failure coverage above.

\paragraph{Worked instance: EnsureMember.}
The mechanized instance inserts a designated element $x$ into a protected
set. Let $B$ be the pre-state combined with permitted environment
additions, and let $L$ be the set of attempts that linearized. Its
effect predicates are
\[
  \mathsf{Zero}_r\equiv S_{\mathrm{post}}=B,
  \qquad
  \mathsf{One}_r\equiv S_{\mathrm{post}}=B\cup\{x\}.
\]
The contract constrains environment additions, requires each member of
$L$ to correspond to an invocation, and relates the post-state to whether
$L$ is empty. A classified success requires its attempt to be in $L$
and its value to have the specified success identifier; a classified
failure requires its attempt not to be in $L$.
Success then supplies a linearized attempt and the specified result.
For a request-wide failure, the broker proves that every invocation has
a classified failure. Since every linearized attempt must be an
invocation, $L$ is empty and the post-state is $B$. Bound follows by
the two cases for $L$ under the contract's state relation.

\paragraph{Deriving the contract from operational steps.}
The EnsureMember operational model supplies a second proof layer. Its
execution predicate is defined by an initial state and enabled state
transitions. A service-linearization step requires a previously invoked,
not-yet-linearized and not-failed attempt; it inserts $x$ into the
protected set and records that attempt in $L$. Environment steps add
only elements other than $x$. Success and failure deliveries are
enabled according to the attempt's linearization state.
An inductive invariant connects the accumulated physical history, $L$,
environment additions, and the protected set. Its preservation implies
the abstract rely relation above for every finite operational execution.
The WAL composition then requires an operational execution and a legal
WAL execution with matching global-event projections; the rely relation
is derived before applying terminal refinement.
% Source: mechanized/t6_adapter_executable_refinement.rs:
% a1_exec; a1_every_exec_configuration_satisfies_invariant;
% ensure_member_exec_derives_adapter_rely;
% ensure_member_executable_wal_terminal_refines.

Two concrete witnesses establish nonvacuity. A typed-WAL \Commit execution
with $x\notin B$ satisfies $\mathsf{One}_r\land\neg\mathsf{Zero}_r$.
An operational crash/retry execution loses the first attempt's success,
records a failed retry, and returns \Unknown with the same effect relation.

% Source: t6_a1_executable_crash_retry_refines and
% t6_a1_executable_crash_retry_nonvacuity in the same module.

\subsection{Deduplicated Operations}

Each attempt carries the same key, bound to the canonical request
within an adapter namespace. The service atomically associates the key
with the operation and a retained result. The contract requires
request binding, decision consistency between effects and repeated
result retrievals, and retention throughout recovery. A retained success
establishes the specified result and update; a retained failure establishes
$\mathsf{Zero}_r$, even when an earlier local outcome is missing.
Repetition retrieves the same decision without another abstract update.
The wrapper distinguishes this retained decision from transport errors
and credential rejections. The mechanized instance proves the laws under
the consistency and retention contract.

\subsection{Instantiation and Reuse}

The read-only instance establishes valid observations without request
mutation at the adapter, Broker, and WAL boundary; stable results across
environment changes require a snapshot or version guarantee.
The uncontrolled profile instead relies on an at-most-one-invocation
broker policy and returns \Unknown for unresolved attempts. Its effect
interpretation uses adapter laws that account for partial and secondary
effects, as required by Bound.
The idempotent and deduplicated instances additionally provide
protected-service couplings, summarized in Section~\ref{sec:evaluation}.

For adapters satisfying the primitive laws, the generic lemma
derives the common semantic interface:
\[
\begin{aligned}
  &\mathsf{SuccessSound}(A)
   \land\mathsf{FailureSound}(A)\\
  &\qquad\land\mathsf{EffectBounded}(A)
   \Rightarrow\mathsf{AdapterVerified}(A).
\end{aligned}
\]
The derivation is a terminal-case split. Each instance supplies its state
relations, rely condition, and law proofs; the broker supplies the causal
evidence and physical-history premises.

\section{Crash-Recovery Protocol and Refinement}
\label{sec:proof}
\label{sec:architecture}

The proof connects retained records to physical events, preserves this
connection through storage refinement, and applies the service laws to
derive an effect claim for each admitted external run. Figure~\ref{fig:a2e-composition}
summarizes the roles of evidence, service laws, and the A2E criterion.

\begin{figure}[t]
  \centering
  \includegraphics[width=\linewidth]{figures/proof-dependencies-v2.pdf}
  \caption{Proof dependencies. The broker derives physical evidence from
  retained records; operational service invariants discharge the adapter
  laws. Their composition yields the per-execution effect theorem. The
  protected-service transport additionally requires coupling and a closed
  interface. A2E lifts the result under uniform observation-class premises.}
  \label{fig:a2e-composition}
  \Description{Retained evidence and operational service laws compose into terminal-effect refinement, with additional protected-service and observation-class premises.}
\end{figure}

\subsection{Protocol Establishes Evidence}

Admission fixes the request identity, canonical-call digest, capability,
budget, recovery profile, and optional key. A request's acknowledged
\texttt{Authorize}, \texttt{Prepare}, and \texttt{Arm} establish its durable
admission and prepared state. Each physical call additionally requires an
acknowledged attempt-specific \texttt{Start}, whose log position identifies
that attempt. Each invocation thus has a unique acknowledged Start and
authorization ancestry; the converse is false because a crash after
Start can prevent invocation.

A response is accepted only for the request and attempt occupying the
executor slot. Each Start has at most one typed Outcome. Terminal
records use strictly backward, request-local references: \Commit cites
the exact successful outcome and value, \Fail cites failure evidence,
and \Unknown records a reason and anchor. Journal legality permits
at most one terminal per request.

Replay reconstructs phase, attempts, budget, revocation status, outcomes,
and terminal state from the acknowledged prefix. Unfinished requests
use retained evidence and the profile-specific classifier; acknowledged
recovery updates survive further crashes, and established terminals are
preserved. Recovery may append a \Commit whose exact successful Outcome
and value were already durable at the crash. Accordingly, committed history
extends monotonically; exact history equality holds for recovery episodes
that append no new \Commit. Induction over ordinary
and recovery transitions connects these records to executor state,
physical invocations and deliveries, budget consumption, and decisions.

\subsection{Refinement Preserves Decisions}

The Broker specifies execution and evidence obligations; the Journal
realizes them with atomic typed records; the typed WAL implements the
Journal with full writes, torn tails, flush cuts, crashes, and truncation:
\[
  \text{typed WAL}\preceq\text{Journal}\preceq\text{Broker}.
\]
Here $\preceq$ preserves the observations used by the contract. Journal
legality enforces record order, references, and request consistency.
WAL refinement preserves the acknowledged prefix without manufacturing
outcomes during tail repair. Monotone prefix mappings allow concrete
steps either to advance the abstract execution or to leave it unchanged.
The simulations preserve request fields, physical-event correspondence,
evidence references, and replay state, so corresponding executions select
the same terminal. Contextual replacement lifts this result to
storage-parametric contexts satisfying the modeled interface conditions,
preserving designated endpoint observations and the selected terminal.

\subsection{Composition Establishes Effects}

Fix an admitted finite execution $\tau$ with journal $J$, request-local
physical history $\eta=\eta_r(\tau)$, and selected terminal $d$. For the
A2E corollary, cut $\tau$ immediately after terminal acknowledgement.
The branch guards and supporting evidence are evaluated in the strict
pre-terminal journal prefix $J^{-}$, giving $H=\mathsf{Obs}_r(J^{-})$.

The composition uses the proof responsibilities in
Table~\ref{tab:proof-responsibilities}. Success and Failure interpret the
selected evidence; Bound also constrains unresolved runs.

\paragraph{Protected-service transport and quantification.}
The service coupling relates calls, returns, and state transitions
at matching prefixes. Coverage accounts for the protected calls, the closed
interface confines calls to the broker, and the rely condition governs
environment steps. The contextual theorem (T6-X0) applies the effect result to supplied
executions, prefix maps, a context containing the WAL, the P0 prefix product,
and a recorded terminal. The EnsureMember witness discharges
these premises for one concrete execution. The shared-WAL family theorem
(DD5) requires a valid operational/protected family member for every terminal
request; the DD2 instance supplies this coverage for a particular execution.
Applying this interface to a production service requires constructing its
matching executions and establishing the same coupling and coverage premises.
ReadOnly is instantiated through the adapter, Broker, and WAL layers
(Table~\ref{tab:proof-surface-main}).

The following theorems summarize the mechanized boundaries.

\begin{theorem}[Provenance and prefix preservation]
\label{thm:history-preservation}
In every admitted finite execution, physical invocations have unique
acknowledged authorization and attempt ancestry; accepted outcomes
retain matching provenance. The request-scoped acknowledged projection
survives crashes and Broker--Journal--typed-WAL refinement.
\end{theorem}

\paragraph{Worked example: why failure coverage spans attempts.}
Return to the deployment request in Figure~\ref{fig:motivating-example}.
Let $a_1$ be the attempt whose reply is lost and $a_2$ the retry rejected
by the service. Write $S_i$ for the acknowledged Start of $a_i$ and $F_i$
for its matching, validated failure Outcome. Table~\ref{tab:worked-evidence}
shows the evidence relevant to the failure guard; authorization records
are implicit.

\begin{table}[t]
  \caption{Evidence for an idempotent request. The first two rows follow
  the motivating execution. The third is a separate history in which the
  first attempt returns a failure accepted under the adapter contract.}
  \label{tab:worked-evidence}
  \centering\small
  \begin{tabularx}{\linewidth}{@{}lY@{}}
    \toprule
    Retained attempt evidence & Supported conclusion \\
    \midrule
    $S_1$ & The first attempt's effect is unresolved; another attempt
      may be eligible. \\
    $S_1,S_2,F_2$ & The retry failed. The first attempt remains unresolved,
      so the classifier selects \Unknown. \\
    $S_1,F_1$ & The sole started attempt has matching failure evidence;
      the failure guard permits \Fail. \\
    \bottomrule
  \end{tabularx}
\end{table}

The guard checks records, while the effect claim concerns physical calls.
Connecting them requires a one-way invariant: every invocation has an
acknowledged Start with authorization ancestry. A Start may survive even
when a crash prevents its call, so record coverage can be conservative.
Outcome provenance supplies the other link: each covering failure record
must identify the corresponding request, attempt, and physical delivery.
Together, these facts turn coverage of Starts into coverage of invocations,
allowing the adapter's Failure law to establish $\mathsf{Zero}_r$.

Crash and replay must preserve both links. Losing an acknowledged Start
could hide an earlier invocation; accepting a failure under the wrong
attempt identity could falsely complete coverage. The Broker invariant
rules out these cases, and the Journal/WAL simulations preserve the
record and physical-history projections used by the argument. This is how
the evidence rule continues to justify request-level outcomes after recovery.

\begin{theorem}[Request-wide terminal evidence after recovery]
\label{thm:terminal-evidence}
Let a well-formed configuration and an admitted Broker execution select
terminal $d$ for $r$, with exact journal and physical-history projections
$J$ and $\eta$. The Broker invariant establishes:
\begin{itemize}
  \item A \Commit cites a matching successful delivery for the exact
    request, attempt, and returned value.
  \item For an idempotent request selecting \Fail, every invoked
    attempt has a matching failure delivery:
    \[
      \forall a\in\mathsf{Invoked}_r(\eta).\quad
      \mathsf{Delivery}(\eta,r,a)=\mathsf{Failure}.
    \]
  \item A \Fail cites its selected failure delivery, and an
    \Unknown reason has the required durable anchor.
\end{itemize}
These causal conclusions establish $\mathsf{OutcomeEvidence}(J,r,\eta,d)$.
Under the adapter rely condition the execution also satisfies
$\mathsf{BrokerOutcomeCompatible}(J,r,\eta,d)$, including
retained-decision consistency for deduplicated operations. The conclusions
hold for corresponding Journal and typed-WAL executions under their
trace-agreement and representation premises.
\end{theorem}
\begin{proof}[Proof outline]
Terminal references resolve to matching deliveries. In the idempotent
failure case, each physical invocation has a durable attempt ancestor;
the failure guard covers every started attempt, and outcome causality
turns this record coverage into invocation-wide failure coverage. Thus
a lost earlier outcome cannot be hidden by a later failure. Deduplicated
compatibility additionally uses decision consistency, while uncontrolled
compatibility uses the invocation bound. Storage refinement supplies
the exact projections and Broker invariant. These are the T6-E0,
T6-C0, and T6-S0 proof boundaries.
\end{proof}
% Evidence: t6_terminal_compatibility.rs,
% idempotent_failure_covers_all_invocations; event_terminal_outcome_is_compatible.

\paragraph{Separating recovery evidence from service semantics.}
Theorem~\ref{thm:terminal-evidence} provides the interface between the
broker's records and the service's effects: causal evidence identifies
matching physical deliveries, and profile compatibility supplies the
premises of the selected service law. For an idempotent request, those
premises include failure coverage over every invocation; for a
deduplicated request, they use the service's retained-decision consistency.
The next theorem composes this evidence with Success, Failure, and Bound
to justify the request-level outcome.
For another service using an existing recovery profile and the same
broker/storage interface, the service-specific work is to define its
effect and result relations, establish the laws and operational rely
invariant, and supply matching trace/representation premises. The broker's
ancestry, outcome-provenance, and storage-preservation proofs remain
reusable; protected-service transport additionally requires coupling and
coverage. EnsureMember and the keyed-decision instance demonstrate this
separation by applying the same generic WAL terminal-refinement theorem
to different service semantics (Table~\ref{tab:proof-surface-main}). This
reuse concerns the connection from retained evidence to request-level
effects, including invocations whose replies were lost.

\begin{theorem}[Per-execution recovery-to-effect refinement]
\label{thm:authorization-to-effect}
Let $C$ be well formed, and let $\tau$ be an admitted finite execution
of the serialized Broker, Journal, or typed-WAL semantics with terminal
$d$ for $r$. At the storage boundaries, require the modeled trace
agreement and representation premises. Let an external run satisfy
the adapter rely condition with respect to $\eta_r(\tau)$, and suppose
the adapter satisfies Success, Failure, and Bound.

Then the selected terminal has the following external interpretation:
\begin{itemize}
  \item \Commit returns the value of its matching successful delivery,
    satisfies $\mathsf{resultSpec}_r$, and establishes $\mathsf{One}_r$
    for a mutating request or $\mathsf{Zero}_r$ for a read-only request.
  \item \Fail has a matching failure delivery and establishes
    $\mathsf{Zero}_r$ for the entire request.
  \item \Unknown makes no conclusive effect claim; the separate Bound
    law still establishes $\mathsf{Zero}_r\lor\mathsf{One}_r$.
\end{itemize}
The history obligations of Theorem~\ref{thm:history-preservation}
continue to hold. A covered operational/protected-service coupling
under the closed interface additionally transports the effect
interpretation to that modeled service execution.
\end{theorem}
\begin{proof}[Proof outline]
Theorem~\ref{thm:terminal-evidence} supplies causal evidence and
profile compatibility for the selected terminal. The adapter's
primitive laws then establish the appropriate result and state
relation for this external run. In the mechanization, their case split
derives $\mathsf{AdapterVerified}$, and the T6-A0 closure applies that
interface to obtain $\mathsf{Refines}$ and per-request effect refinement.
The Journal/WAL wrappers retain the same evidence and terminal;
the covered service coupling transports the state relation at the
additional protected-service boundary.
\end{proof}

\begin{corollary}[A2E terminal soundness]
\label{cor:a2e-soundness}
Fix $C$, $A$, $r$, and an observation $H$ realized by an execution in
$\mathsf{DecExec}_{C,A}(r)$. Suppose the Armed-phase classifier selects
terminal $d$ from $H$, where $d$ records the terminal kind and the
returned value for \Commit. Evidence references and reason annotations
must separately satisfy the terminal-append rules. Require the
operational, projection, and adapter
premises of Theorem~\ref{thm:authorization-to-effect} for every
execution in $\mathsf{Worlds}_{C,A}(H,r)$.
Then
\[
 \mathsf{Compat}_{C,A}(H,r)\subseteq\llbracket d\rrbracket_{r,A},
 \qquad H,r,C,A\vdash_{\mathrm{A2E}}d.
\]
\end{corollary}
\begin{proof}
For any $\tau\in\mathsf{Worlds}_{C,A}(H,r)$, terminal legality
checks the classifier against the same $H$. Fixed $C$ and $r$ therefore
give the same terminal kind and, for \Commit, returned value. Apply
Theorem~\ref{thm:authorization-to-effect} to each such terminal-containing
execution and collect its external outcomes. For \Unknown the inclusion
is immediate; Bound is a separate conclusion.
\end{proof}
The corollary is a paper-level lift of the mechanized per-execution theorem;
its premises must hold throughout the observation class.
Table~\ref{tab:recovery-witnesses} gives concrete mechanized executions
satisfying the premises, including successful and unresolved outcomes.

\begin{table}[t]
\caption{Mechanized recovery witnesses and the evidence they establish.
These are concrete executions, separate from the dynamic crash campaign.}
\label{tab:recovery-witnesses}
\small\centering
\begin{tabularx}{\linewidth}{@{}lYY@{}}
\toprule
Witness & Retained evidence and service behavior & Justified result \\
\midrule
Durable success (H2) & Successful Outcome survives the crash; terminal is absent &
Recovery appends \Commit with that value; committed history strictly extends \\
EnsureMember (A1/X0) & First attempt changes the set but its success is not retained;
retry fails & \Unknown retains the unresolved earlier effect \\
Keyed decision (DD2--DD5) & First attempt applies and memoizes; retry retrieves
the same successful decision & \Commit with one abstract update under the key contract \\
\bottomrule
\end{tabularx}
\end{table}

\section{Executable Realization}
\label{sec:implementation}

A verified typed-record kernel checks candidate records and recovery
transitions on selected Rust broker paths. The Rust implementation supplies
request handling, byte persistence, and service invocation.

\subsection{Reference Broker and Verified Kernel}

The reference broker is implemented in Rust using only the
standard library. It allocates monotone request identifiers,
enforces capability budgets, maintains one executor slot, and
stores records in a checksummed byte WAL with torn-tail recovery.
Example adapters exercise idempotent, deduplicated, and
uncontrolled mutating operations. The reference broker rejects \ReadOnly requests; that instance is covered by
formal proofs and witnesses. In the deduplicated harness, service-owned
decisions survive broker restarts.

The verified kernel operates on the typed alphabet \texttt{Authorize},
\texttt{Prepare}, \texttt{Arm}, \texttt{Start}, \texttt{Outcome}, and the
terminal records \texttt{Commit}, \texttt{Fail}, and \texttt{Unknown}.
It maintains replay-derived request state, the Journal projection, log
sequence number, and acknowledgement cut. References are request-local
and strictly backward.

Kernel verification connects executable guards and accepted state updates
to abstract record legality and replay. It preserves request ownership and
backward references, phase and budget agreement, exact recovery \Commit
provenance, reason-specific \Unknown anchors, and crash/recovery control.
These properties hold for legal sequences over the supported typed alphabet.
The kernel rejects \texttt{Revoke}.

\paragraph{Executable classification.}
Rust's \texttt{recovery\_decision} follows the guard order in
Section~\ref{sec:contract}. An Idempotent Failure with incomplete coverage
returns \Unknown(NonConclusiveFailure). Deduplicated Failure follows its
keyed-decision contract. Ambiguous and InvalidResult permit retry for
supported repetition-safe profiles when attempts remain. An Uncontrolled
request stops after a recorded outcome or pending Start.

\subsection{Gated Append and Recovery}

\paragraph{Append.}
Before writing a candidate record, the gate asks the kernel to
preview it against the current checked state. Only accepted candidates proceed to the gated WAL write. After the WAL
returns its durable log sequence number, the integration commits
the corresponding transition to the kernel:
\[
  \text{preview}
  \;\longrightarrow\;
  \text{WAL write}
  \;\longrightarrow\;
  \text{kernel commit}.
\]
The final step updates kernel state after acknowledged persistence; it is
separate from the request's \Commit terminal.

\paragraph{Reopen and replay.}
Before a reopened broker serves requests, the gate replays the
full acknowledged prefix through the kernel. This reconstructs
the checked request state from retained records, including
terminal evidence and remaining budgets. Recovery transitions
are then subject to the kernel's recovery guards and
class-sensitive resume predicate. For an Armed request, the kernel permits
Idempotent and Deduplicated resumption; Uncontrolled and ReadOnly resumption
requires that no durable Start exists. The kernel supports ReadOnly resumption under this guard, although the
reference broker rejects that profile at admission.

\paragraph{Fail-closed operation.}
The fail-closed evaluation profile requires an immutable manifest
and a gate for every broker open. Missing required checks or
rejected transitions stop the gated workflow. The harness records gate/open coverage,
preview-to-commit pairing, full-prefix replay, and recovery
completion, making the use of the gate observable in the
evaluation traces.

\subsection{Verification and Integration Scope}

Kernel verification establishes invariant preservation for supported typed
transitions. The integration campaign checks trace conformance on selected
Rust paths. Record translation, gate routing, byte decoding, checksums,
filesystem durability, and surrounding Rust control flow form the trusted
implementation boundary. Services must satisfy their adapter laws.

\section{Evaluation}
\label{sec:evaluation}

RQ1 tests recovery trace conformance and the failure-coverage rule under
broker crashes. RQ2 audits common proof interfaces and service-specific
obligations. RQ3 measures serialized local execution cost.

\subsection{RQ1: Do Recovery Paths Preserve Evidence and Make Sound
Effect Decisions?}

\paragraph{Trace conformance.}
The fail-closed campaign checks append and recovery transitions through
the verified typed-record kernel. Its gate compares retained records,
acknowledgement cuts, evidence references, and replay-derived decision
state and budgets. Each crash case also asserts terminal uniqueness,
invocation authorization ancestry, the terminal attempt limit, the
expected terminal decision, and the example adapter's effect or
mutation count, observed independently of the kernel state comparison.

The retained gated run records 47 broker opens and 47 gates, 1,573
preview/commit pairs, and 675 replayed records. Of 47 recovery instances,
40 finish recovery and seven take a verified resume path.

The matrix exercises three mutating classes at seven crash sites:
after authorization, preparation, arming, attempt start, invocation,
outcome persistence, and terminal persistence. All 21 cases satisfy
these assertions: 19 end in \Commit and two in \Unknown, with no \Fail.
The two \Unknown cases are uncontrolled requests interrupted after
Start and after invocation, with observed effect counts zero and one,
respectively. Separate stale-delivery tests assert that rejected deliveries leave the WAL and
executor slot unchanged. For the synchronous uncontrolled adapter,
the send-to-linearization window is represented in the model; the
harness observes only the boundaries surrounding the invocation.

\paragraph{Effect discrimination.}
To test the request-wide failure-coverage rule after a real process
crash, we constructed paired broker executions with indistinguishable
durable evidence but different external effects. An independently
running Redis 7.2.5 service executed the idempotent \texttt{SADD}.
In the zero-effect execution, the broker was killed with \texttt{SIGKILL}
after Start but before sending the command. In the one-effect execution,
a proxy let Redis apply \texttt{SADD} but withheld its reply before the
broker was killed. Redis remained alive in both. After broker restart,
an ACL change revoked \texttt{SADD}, so the retry returned \texttt{NOPERM}.
The target member is initially absent in the isolated service state.
An independent \texttt{SISMEMBER} oracle observes absence in the first
execution and presence in the second.

Within each classifier variant, the paired executions have byte-identical
initial and pre-terminal WALs. The complete classifier returns \Unknown
in both; removing only the Idempotent incomplete-failure-coverage
branch returns \Fail in both. Table~\ref{tab:redis-crash} also shows
that conclusive-failure and success controls retain their decisions.
Since \Fail asserts zero request effect, the ablated one-effect execution
violates the decision contract. Reopening terminal requests made zero additional service calls.

\begin{table*}[t]
  \caption{Redis broker process-crash campaign. Within each classifier
  variant, the first two rows have byte-identical initial and
  pre-terminal WALs, despite different service effects.}
  \label{tab:redis-crash}
  \centering
  \small
  \begin{tabular}{lccc}
    \toprule
    Scenario & Oracle membership & Full classifier & Coverage ablation \\
    \midrule
    Paired: before send & Absent & \Unknown & \Fail \\
    Paired: reply lost & Present & \Unknown & \Fail \\
    Conclusive failure & Absent & \Fail & \Fail \\
    Success & Present & \Commit & \Commit \\
    \bottomrule
  \end{tabular}
\end{table*}

\paragraph{What the two campaigns establish.}
The Redis ablation demonstrates the consequence of omitting request-wide
failure coverage: identical retained evidence leads to an unsound \Fail
in the execution with an effect. The success and conclusive-failure controls
show that sufficient evidence still permits decisive outcomes. The kernel
campaign checks the append and recovery paths that preserve this evidence
in the exercised Rust executions.

\subsection{RQ2: Which Proof Interfaces Do the Service Instances Reuse?}

Table~\ref{tab:proof-surface-main} traces each instance to its common
proof interfaces, service-specific obligations, and demonstrated boundary. The
EnsureMember and keyed-decision proofs both call the same generic
WAL terminal-refinement theorem after establishing their adapter laws,
rely relation, and trace/representation premises. The read-only witness
uses the common WAL/Broker representation and semantic definitions but
establishes its concrete effect refinement directly.
% Audit anchors: t6_adapter_executable_refinement.rs:
% ensure_member_executable_wal_terminal_refines -> wal_terminal_outcome_refines;
% t6_deduplicated_witness.rs: dd2_wal_terminal_refines ->
% wal_terminal_outcome_refines; t6_readonly_operational.rs:
% ro_wal_closed_representation, ro_concrete_refines, ro_wal_terminal_refines.

\begin{table*}[!t]
  \caption{Structural proof reuse: common theorem dependencies,
  service-specific obligations, and demonstrated scope.}
  \label{tab:proof-surface-main}
    \centering
    \small
    \setlength{\tabcolsep}{4pt}
    \begin{tabularx}{\textwidth}{@{}lYYY@{}}
      \toprule
      \textbf{Instance} & \textbf{Common proof interface used}
      & \textbf{Instance obligations} & \textbf{Demonstrated boundary} \\
      \midrule
      EnsureMember
      & Generic WAL terminal refinement, using common terminal evidence
        and profile compatibility
      & Set-update semantics, response classification, operational rely
        invariant, and service coupling
      & Operational/WAL refinement plus covered protected-service and
        contextual construction \\
      Read-only sample
      & Common WAL/Broker representation and effect definitions;
        direct concrete effect proof
      & Observation validity under environment changes, zero request
        effect, and operational rely invariant
      & Concrete adapter/WAL/Broker witness; no protected-service or
        contextual instantiation \\
      Keyed decision
      & Same generic WAL terminal-refinement theorem as EnsureMember
      & Key/request binding, retained-decision consistency, operational
        rely invariant, and coverage
      & Covered protected-service/contextual witness and
        coverage-conditioned request family \\
      \bottomrule
    \end{tabularx}
\end{table*}

This dependency audit measures structural reuse. Effort savings require
additional controls and person-hour measurements.
The historical 56-target proof-effort inventory is not combined with
the later verification campaign's counts.
% Historical inventory: reference-broker/evaluation/results/rq3-proof-effort.json.

\subsection{RQ3: What Is the Local Cost of the Selected Serialized
Implementation?}

\begin{table}[t]
    \caption{Serialized local cost: mediated and journaled workloads.}
    \label{tab:serialized-local-cost}
    \centering
    \scriptsize
    \setlength{\tabcolsep}{2pt}
    \begin{tabularx}{\linewidth}{@{}Yrrrr@{}}
      \toprule
      \textbf{Mode} & \textbf{ms/req} & \textbf{95\% CI} &
      \textbf{req/s} & \textbf{Rec. ms} \\
      \midrule
      Mediated & 0.464 & 0.452--0.475 & 2,165.87 & 1.268 \\
      Journaled & 0.144 & 0.139--0.148 & 7,011.15 & 0.020 \\
      \bottomrule
    \end{tabularx}
\end{table}

Table~\ref{tab:serialized-local-cost} compares the mediated
broker with a journaled at-least-once ablation. Recovery time measures
reopening and replaying the WAL after the primary workload. A separate
in-memory direct loop reaches $13.9\times10^6$ requests/s; its different
workload supplies a latency floor. The journaled ablation provides the
workload comparison.

The mediated workload costs 0.320~ms more per request than the
journaled ablation in this campaign, with 3.22$\times$ its mean latency.
The workloads differ in record and flush counts; the measured difference
is the end-to-end cost of the selected design.

Ambiguous retry workloads average 0.588~ms for \Idempotent,
0.592~ms for \Deduplicated, and 0.285~ms for the at-least-once
ablation. A separate retained storage campaign reports 0.391~ms on
server overlayfs, 0.0358~ms on server tmpfs, 0.456~ms on server
NVMe-backed ext4, and 17.174~ms on an independent WSL2 ext4 virtual
disk. The variation shows that the absolute cost depends strongly
on the storage path.

The measurements characterize serialized local execution with the stated
durability policy. Equal-flush batching, concurrency, remote-tool latency,
and recovery scaling with log size remain unmeasured.

\subsection{Reproducibility}

The 21-case campaign uses one serialized executor and the same named
fault sites for each profile. Separate recovery-interruption tests
reopen after Start, interrupt again after Outcome, and reopen once
more. The Redis runner builds the full and single-branch ablation variants
and retains tool versions, source hashes, WALs, service-oracle values, and
per-case logs. Redis ran with AOF and \texttt{appendfsync always}.
The reproduction command and scope are documented in
\texttt{reference-broker/evaluation/redis-process-crash.md}.
The earlier four-scenario reference-service run is retained as
\texttt{decision-discrimination.json}. A separate formal campaign reports
72 passing registered Verus targets. The report records 1,534
dependency-aware nonduplicated obligations; the sum across targets is
39,706 because cumulative targets recheck dependencies. The recorded
Verus version is \texttt{0.2026.07.05.49b8806}, with Rust 1.96.0.
The retained source-policy scan passes its proof-escape checks.
The bounded TLA+ full suite reports 13 of 13 passing scenarios using
Java 21.0.11 and TLC 1.7.4.

The benchmark uses five warmups and 30 measured repetitions, each with
500 primary requests and 50 retry requests on one logical processor and
local overlayfs. Every repetition runs in a fresh process with a fresh
temporary WAL. The 95\% normal approximation intervals use independent
run-level measurements. Throughput averages each repetition's aggregate
wall-clock rate rather than inverting mean per-request latency. Source,
broker, kernel, Verus, and Rust toolchain hashes identify inputs to the
individual retained campaigns. These reports describe separate source
snapshots: K4-I1 names revision \texttt{1e37979},
the TLA+ full report names \texttt{9a45d39}, and the Verus report identifies
its isolated source copy. Performance and Redis reports retain their own
provenance. Source/report reconciliation with the current development tree
remains in progress. The historical 18-cell
fault-window matrix is a conceptual coverage map, including model-only
windows, rather than an additional executed campaign.

\section{Discussion and Limitations}
\label{sec:limitations}

\paragraph{Deployment and semantic scope.}
Protected-service transport requires a closed interface with confined
handles, authenticated endpoints, correct request/adapter binding, and
prevention of bypass. Deployment must establish these conditions and the
adapter's rely and retention guarantees. For EnsureMember, environment
steps preserve the target member. Competing rollbacks or overwrites require
stronger ownership/version guarantees or a relation recording historical
application. An in-flight invocation may take effect after restart,
including after the broker returns \Unknown.

\paragraph{Safety and implementation boundary.}
The results cover finite serialized fail-stop histories under trusted
configuration and typed-WAL semantics. Open obligations include liveness,
fairness, concurrent linearizability, planner delivery, production-adapter
conformance, and refinement of the implementation boundary in
Section~\ref{sec:implementation}. The Redis experiment tests broker crashes
with a live service; Redis correctness and power-loss durability require
separate evidence.

\paragraph{Application scope.}
Agents motivate request synthesis and immediate planner dependence on tool
outcomes. Other orchestrators can use the same contract when they provide
stable request identity, retained evidence, and appropriate service laws.

\section{Related Work}
\label{sec:related}

\paragraph{Authorization and governed agents.}
Capability systems establish authority and confinement boundaries
\cite{saltzer1975protection,lampson1974protection,watson2010capsicum}.
AgentBound applies resource capabilities to MCP-server execution
\cite{buehler2026agentbound}; CaMeL separates trusted control flow from
untrusted data with policies at tool calls~\cite{debenedetti2025camel}.
Guardians proposes verification of structured agent workflows before
execution~\cite{meijer2025guardians}. Our recovery model addresses the
temporal question left when an authorized invocation's result is lost:
which request-level effect claim is supported after restart? These
authorization mechanisms can help realize the closed-interface premise
used by our service-coupling theorem.

Proof of Execution (PoE) attests authorization, path compliance, trace
integrity, and replayability under cryptographic and deployment
assumptions~\cite{rhodes2026poe}. Our recovery invariant complements trace
attestation by preserving the unresolved status of lost responses across
later failed attempts.

\paragraph{Durable workflows and service-owned decisions.}
Temporal reconstructs workflow state from durable history while external
activities require appropriate retry semantics
\cite{temporal2026workflow,temporal2026activity}. Reliable State Machines
provides failure transparency for its programming framework
\cite{mukherjee2019rsm}, and Beldi provides fault-tolerant transactional
stateful serverless workflows~\cite{zhang2020beldi}. RIFL uses retained
completion records and request identifiers to provide linearizable
exactly-once operations~\cite{lee2015rifl}; idempotency-key interfaces make
service-owned evidence available under their binding and retention contracts
\cite{jena2025idempotency,stripe2026idempotency}. Our adapter interface states which service guarantees justify a broker
terminal, including the distinction between retrieving a retained decision
and repeating an idempotent attempt whose reply is missing.

\paragraph{Machine-checked recovery and refinement.}
FSCQ and Argosy verify crash-safe storage and layered recovery, while
Perennial supports reasoning about concurrency and crash safety
\cite{chen2015fscq,chajed2019argosy,chajed2019perennial}.
Disel, IronFleet, and Verdi provide machine-checked distributed-system
reasoning and refinement~\cite{sergey2018disel,hawblitzel2015ironfleet,
wilcox2015verdi}. seL4 establishes implementation refinement at an
operating-system kernel boundary~\cite{klein2009sel4}.
Our development specializes layered recovery refinement to the connection
between durable broker evidence and physical invocation histories.
This connection supplies the premises of service-adapter laws, yielding
request-level effect claims through a common terminal-refinement interface.
The storage interface is a typed WAL; a verified typed-record kernel
checks selected executable broker paths.

Dual-Write Recovery formalizes delivery across independent source and sink
authorities in Isabelle/HOL~\cite{andreakis2026dualwrite}. It studies
source-only information bounds, authoritative sink evidence, in-flight
sends, concurrent recoverers, and evidence lifetime. A2E's theorem interface
connects an authorizing broker's durable records to invocation-wide
terminal evidence, then composes that evidence with service-specific
effect laws. The broker/storage argument is shared by the EnsureMember
and keyed-decision instances. Our serialized model assumes retention
throughout recovery.

\section{Conclusion}
\label{sec:conclusion}

We verify request-level outcome reporting for a serialized recovering
broker. A durable evidence invariant connects authorization and failure
records to physical invocations; storage refinement preserves it, and
adapter laws give it effect meaning. A verified kernel checks selected
Rust paths. Paired Redis crashes demonstrate why request-wide failure
coverage matters: removing it produces an unsound \Fail after a lost success.

\section*{Data Availability}
The project contains the formal models, Verus sources, reference broker,
experiment runners, and retained verification and evaluation reports.
An anonymous reviewer-accessible release has not yet been made available,
so this draft provides no artifact download link. The principal theorem
premises, experimental procedures, and reported results are described in
the paper.
% BEFORE SUBMISSION: replace this draft status with the real anonymous
% artifact location and public-release intention. Do not link an
% author-identifying repository in the anonymous manuscript.

\bibliographystyle{ACM-Reference-Format}
\bibliography{references-ndss-final}

\end{document}
